\documentclass{article}
\usepackage{amsmath,amssymb}
\usepackage{xurl}
\usepackage[hidelinks]{hyperref}
% Shared commands for notes in docs/paper-gaps/.

\DeclareUnicodeCharacter{2080}{\ifmmode{}_{0}\else\textsubscript{0}\fi}
\DeclareUnicodeCharacter{2081}{\ifmmode{}_{1}\else\textsubscript{1}\fi}
\DeclareUnicodeCharacter{2082}{\ifmmode{}_{2}\else\textsubscript{2}\fi}
\DeclareUnicodeCharacter{2099}{\ifmmode{}_{n}\else\textsubscript{n}\fi}

\newcommand{\Lean}{\textsc{Lean}}
\ExplSyntaxOn
\NewDocumentCommand{\leanid}{m}{
  \ifmmode
    \textsf{\detokenize{#1}}
  \else
    \group_begin:
    \urlstyle{sf}
    \tl_set:Nn \l_tmpa_tl {#1}
    \tl_replace_all:Nnn \l_tmpa_tl {\_} {_}
    \exp_args:NV \path \l_tmpa_tl
    \group_end:
  \fi
}
\ExplSyntaxOff
\providecommand{\tr}{\operatorname{tr}}
\newcommand{\tnleanIssueBase}{https://github.com/LionSR/TNLean/issues/}
\newcommand{\tnleanPrBase}{https://github.com/LionSR/TNLean/pull/}
\newcommand{\tnleanDocBase}{https://sirui-lu.com/TNLean/docs/}
\newcommand{\ghissue}[1]{\href{\tnleanIssueBase#1}{GitHub issue \##1}}
\newcommand{\ghpr}[1]{\href{\tnleanPrBase#1}{PR \##1}}
\newcommand{\doclink}[1]{\href{\tnleanDocBase#1.html}{\path{#1}}}

% Dirac notation and the identity operator, as in blueprint/src/macros/common.tex.
% \providecommand keeps note-local definitions authoritative.
\usepackage{dsfont}
\providecommand{\ket}[1]{|#1\rangle}
\providecommand{\bra}[1]{\langle #1|}
\providecommand{\braket}[2]{\langle #1 | #2 \rangle}
\providecommand{\ketbra}[2]{|#1\rangle\!\langle #2|}
\providecommand{\Id}{\mathds{1}}
\providecommand{\one}{\mathds{1}}
\providecommand{\C}{\mathbb{C}}
\providecommand{\MN}[1]{M_{#1}(\C)}
\providecommand{\MD}{M_D(\C)}

% External referencing for paper sources.  The build helper writes sanitized
% auxiliary files under build/paper-aux/ and excludes duplicated source labels.
\usepackage{xr}

% Import a generated paper auxiliary file with a caller-supplied label prefix.
% The candidate paths support builds from docs/paper-gaps/, the repository root,
% and build/paper-aux/ itself.
\newcommand{\importpaperaux}[2]{%
  \IfFileExists{../../build/paper-aux/#2.aux}{%
    \externaldocument[#1]{../../build/paper-aux/#2}%
  }{%
    \IfFileExists{build/paper-aux/#2.aux}{%
      \externaldocument[#1]{build/paper-aux/#2}%
    }{%
      \IfFileExists{#2.aux}{%
        \externaldocument[#1]{#2}%
      }{}%
    }%
  }%
}

% General external reference macros
\newcommand{\paperref}[2]{\ref{#1-#2}}
\newcommand{\papereqref}[2]{(\ref{#1-#2})}

% CPSV16 source (1606.00608 / MPDO-22-12-17-2.tex) external document and shortcuts
\importpaperaux{cpsv16-}{1606.00608/MPDO-22-12-17-2-tnlean}
\newcommand{\cpsvref}[1]{\paperref{cpsv16}{#1}}
\newcommand{\cpsveqref}[1]{\papereqref{cpsv16}{#1}}

% Verdict marker: \gapnote{<kind>}{<status>}. Kind is the policy
% classification (clarification, local-correction, scope-restriction,
% unfaithful, false-source, open-gap); status is open, wip (elimination
% underway), resolved, or historical. Machine-read by the site index;
% prints a verdict line.
\newcommand{\gapnote}[2]{\par\noindent\textbf{Verdict:} #1 (#2).\par}

\title{The Fixed-Point State of a Tensor with Repeated Blocks}
\author{}
\date{2026-09-26}
\begin{document}
\maketitle
\gapnote{false-source}{open}

\section{What the source states}
Let $A^i=\bigoplus_{j=1}^b\operatorname{diag}(\mu_{j,1},\dots,\mu_{j,m_j})
    \otimes A_j^i$ with normal $A_j$ forming a basis of normal tensors
\cite[eq.~(S2)]{Malz2024Preparation}, block $q$ sites, and write the polar
decomposition of the blocked map as $B=VP$. The source asserts that $P$
inherits the block structure of $A$,
\begin{align}
    P & =\bigoplus_{j=1}^b\operatorname{diag}(\mu_{j,1}^q,\dots,\mu_{j,m_j}^q)
    \otimes P_j
    \label{eq:mswc24_block_form}
\end{align}
with normal $P_j$ \cite[eq.~(S5)]{Malz2024Preparation}, replaces each $P_j$ by
its fixed point, and obtains the approximating state
\begin{align}
    \ket{\widetilde\phi_N}
     & \propto V^{\otimes M}\sum_{j=1}^b\beta_j\ket{\Omega_j},
    \qquad \beta_j=\sum_{k=1}^{m_j}\mu_{j,k}^N,\qquad
    \ket{\Omega_j}=\bigotimes_{k=1}^M\ket{\omega_j}_{R_kL_{k+1}},
    \label{eq:mswc24_approx_state}
\end{align}
with $N=qM$ \cite[eq.~(S7)]{Malz2024Preparation}. Part~(ii) of the
approximation-error lemma of the Supplemental Material then bounds
$\varepsilon=1-\lvert\braket{\widetilde\phi_N}{\phi_N}\rvert$ by
$O\bigl((N/q)e^{-\gamma q/\xi_{\mathrm{diag}}}\bigr)$ for $q=o(N)$, with
$\xi_{\mathrm{diag}}$ the largest correlation length of the blocks.

\section{The counterexample}
Take $d=2$ and two normal blocks of bond dimension one, $A_1=(1,0)$ with
$m_1=2$ and $\mu_1=(1,1)$, and $A_2=(0,1)$ with $m_2=1$ and $\mu_2=1$. Then
$A^0=\operatorname{diag}(1,1,0)$, $A^1=\operatorname{diag}(0,0,1)$, $D=3$, and
\begin{align}
    \ket{\phi_N}\propto2\ket{0\cdots0}+\ket{1\cdots1},\qquad
    \beta=(2,1).
    \label{eq:mswc24_counterexample_target}
\end{align}
For every $q$ the only nonzero blocked matrices are
$B^{0\cdots0}=\operatorname{diag}(1,1,0)$ and $B^{1\cdots1}=\operatorname{diag}(0,0,1)$.
As vectors of $\C^D\otimes\C^D$ they are $e_{00}+e_{11}$ and $e_{22}$, so
$P$ has rank two with range spanned by $(e_{00}+e_{11})/\sqrt2$ and $e_{22}$, and
$V$ maps these to $\ket{0\cdots0}$ and $\ket{1\cdots1}$. On the copy index of the
first block, $P$ is the rank-one matrix
$\tfrac1{\sqrt2}\bigl(\begin{smallmatrix}1&1\\1&1\end{smallmatrix}\bigr)$,
the square root of the restriction
$\bigl(\begin{smallmatrix}1&1\\1&1\end{smallmatrix}\bigr)$ of $B^\dagger B$,
rather than
$\operatorname{diag}(\mu_{1,1}^q,\mu_{1,2}^q)=\mathbf 1$ times a positive part, so
(\ref{eq:mswc24_block_form}) fails.

The blocks are one-dimensional, so every block has the trivial fixed point,
$\omega_1$ is $e_{cc}$ for a copy $c\in\{0,1\}$ and $\omega_2=e_{22}$. Then
$Ve_{cc}=\ket{0\cdots0}/\sqrt2$ and $Ve_{22}=\ket{1\cdots1}$, so
\begin{align}
    V^{\otimes M}\sum_j\beta_j\ket{\Omega_j}
    =2\cdot2^{-M/2}\ket{0\cdots0}^{\otimes M}+\ket{1\cdots1}^{\otimes M},
    \label{eq:mswc24_counterexample_approx}
\end{align}
whose overlap with the normalized target (\ref{eq:mswc24_counterexample_target}) is
$(2^{2-M/2}+1)/\bigl(\sqrt5\,(2^{2-M}+1)^{1/2}\bigr)$. It equals $0.988$, $0.949$,
$0.800$ and $0.555$ at $M=1,2,4,8$, tends to $1/\sqrt5$, and does not depend on
$q$. Taking for instance $q\approx\sqrt N$, so that $q=o(N)$ and $M\to\infty$, the
error tends to $1-1/\sqrt5$, while the printed bound tends to zero.

Placing the pair of the first block on both copies does not help. With
$\omega_1=(e_{00}+e_{11})/\sqrt2$, the two legs at one site come from different
pairs, $V$ annihilates the configurations whose two legs lie in different copies,
and only the two copy-uniform configurations survive, each with amplitude
$2^{-M}$, so the first term is suppressed by $2^{1-M}$ relative to the
target (\ref{eq:mswc24_counterexample_target}), instead of the factor
$2^{-M/2}$ of (\ref{eq:mswc24_counterexample_approx}). The failure is not in
the choice of pair: the fixed point of the positive part of the first block is not a product of nearest-neighbour pairs
carrying the weight $\beta_1$.

\section{What is formalized}
The fixed-point pairs of the blocks of a canonical form, each placed on one copy
of its block, the state $\sum_j\alpha^{(N)}_j\ket{\Omega_j}$ and its GHZ form are
formalized.\footnote{Formal declarations:
    \leanid{MPSTensor.SectorDecomposition.embeddedFixedPointPair} and
    \leanid{MPSTensor.SectorDecomposition.canonicalFixedPointState} in
    \doclink{TNLean/MPS/Preparation/NonNormalCanonicalForm}.}
The counterexample is formalized for the source's construction with the pairs
placed on one copy: for every $q\geq1$ and $M\geq1$ the overlap of
(\ref{eq:mswc24_counterexample_approx}) with the normalized target is
$(4\cdot2^{-M/2}+1)/\bigl(\sqrt5\,(4\cdot2^{-M}+1)^{1/2}\bigr)$, which tends to
$1/\sqrt5$. The transfer maps of the two one-dimensional blocks have no eigenvalue
other than one, so every bound $\lvert\lambda_2\rvert<1$ on their subleading
eigenvalues is admissible and $e^{-\gamma/\xi_{\mathrm{diag}}}<1$; along $q=M$,
$N=M^2$, which satisfies $q=o(N)$, the error exceeds
$C(N/q)e^{-\gamma q/\xi_{\mathrm{diag}}}\exp\bigl(C(N/q)e^{-\gamma
q/\xi_{\mathrm{diag}}}\bigr)$ for every constant $C$.\footnote{Formal declarations:
    \leanid{MPSTensor.nonNormalApproxOverlap_repeatedBlockTensor},
    \leanid{MPSTensor.tendsto_norm_nonNormalApproxOverlap_repeatedBlockTensor},
    \leanid{MPSTensor.not_approximationError_le_repeatedBlockTensor} in
    \doclink{TNLean/MPS/Preparation/RepeatedBlockCounterexample}.}
The positive part and the support projector of the blocked tensor are computed
from the uniqueness of the positive square root of $B^\dagger B$. The state with
the pair of the first block spread over both copies is not formalized.

The block form (\ref{eq:mswc24_block_form}) also fails when every multiplicity is
one, as soon as the $q$-site states of distinct blocks are not orthogonal
\cite{gap:mswc24_block_form_mixed_overlap}. Part~(ii) of the lemma therefore has
no formalized version, and restricting to multiplicity one does not give one.

\section{Elimination}
A correct statement for repeated blocks must replace $\sum_j\beta_j\ket{\Omega_j}$
by the fixed point of the actual positive part, whose restriction to the copies of
block $j$ has rank one in the copy index, and carry the weights through $V$. It
must also account for the overlaps of distinct blocks recorded in
\cite{gap:mswc24_block_form_mixed_overlap}. The first step is a corrected form of
(\ref{eq:mswc24_block_form}) for blocks whose $q$-site states are orthogonal.

\bibliographystyle{alpha}
\bibliography{references}
\end{document}
